\section{Exercices d'application}

\rubrique{Nombre dérivé et tangente}

\exo{}\diff{1}
Soit $f$ la fonction définie sur $\R$ par $f(x)=x^2-3x+1$.
\begin{enumerate}[label=\alph*)]
\item Calculer le taux d'accroissement de $f$ entre $1$ et $1+h$ ($h\neq0$).
\item En déduire que $f$ est dérivable en $1$ et donner $f'(1)$.
\item Déterminer l'équation de la tangente à la courbe de $f$ au point d'abscisse $1$.
\end{enumerate}

\exo{}\diff{1}
On considère la fonction $g$ définie sur $\R\setminus\{2\}$ par $g(x)=\dfrac{3x-1}{x-2}$.
\begin{enumerate}[label=\alph*)]
\item Calculer le taux d'accroissement de $g$ entre $3$ et $3+h$ ($h\neq0$).
\item Montrer que $g$ est dérivable en $3$ et déterminer $g'(3)$.
\item Donner une équation de la tangente à la courbe de $g$ au point d'abscisse $3$.
\end{enumerate}

\exo{}\diff{2}
Soit $h$ la fonction définie sur $\R$ par $h(x)=\sqrt{x^2+9}$.
\begin{enumerate}[label=\alph*)]
\item Calculer le taux d'accroissement de $h$ entre $4$ et $4+h$ ($h\neq0$).
\item En déduire que $h$ est dérivable en $4$ et préciser $h'(4)$.
\item Déterminer l'équation de la tangente à la courbe de $h$ au point $A(4;5)$.
\end{enumerate}

\rubrique{Calcul de fonctions dérivées}

\exo{}\diff{1}
Déterminer la fonction dérivée de chacune des fonctions suivantes définies sur $\R$ :
\begin{enumerate}[label=\alph*)]
\item $f(x)=5x^3-2x^2+7x-4$
\item $g(x)=-\dfrac{1}{2}x^4+3x^2-5x+1$
\item $h(x)=2x^5-3x^3+4x-8$
\end{enumerate}

\exo{}\diff{1}
Calculer la dérivée des fonctions suivantes :
\begin{enumerate}[label=\alph*)]
\item $f(x)=\dfrac{2x+1}{x-3}$ sur $\R\setminus\{3\}$
\item $g(x)=(3x-2)(x^2+1)$ sur $\R$
\item $h(x)=\dfrac{x^2-4}{x+2}$ sur $\R\setminus\{-2\}$
\end{enumerate}

\exo{}\diff{2}
Soit $f$ la fonction définie sur $\R$ par $f(x)=(2x-1)^3$.
\begin{enumerate}[label=\alph*)]
\item Développer $f(x)$ puis calculer $f'(x)$.
\item Retrouver $f'(x)$ en utilisant la formule de dérivation d'une puissance.
\end{enumerate}

\exo{}\diff{2}
On considère la fonction $f$ définie sur $\R$ par $f(x)=\dfrac{x^2+1}{x^2+2}$.
\begin{enumerate}[label=\alph*)]
\item Calculer $f'(x)$.
\item Résoudre l'équation $f'(x)=0$.
\end{enumerate}

\rubrique{Dérivée et sens de variation}

\exo{}\diff{1}
Soit $f$ la fonction définie sur $\R$ par $f(x)=x^3-3x+2$.
\begin{enumerate}[label=\alph*)]
\item Calculer $f'(x)$.
\item Étudier le signe de $f'(x)$ et en déduire les variations de $f$.
\item Dresser le tableau de variation de $f$.
\end{enumerate}

\exo{}\diff{2}
Soit $g$ la fonction définie sur $\R\setminus\{1\}$ par $g(x)=\dfrac{x^2-2x+3}{x-1}$.
\begin{enumerate}[label=\alph*)]
\item Calculer $g'(x)$.
\item Étudier le signe de $g'(x)$.
\item En déduire les variations de $g$.
\end{enumerate}

\exo{}\diff{2}
On considère la fonction $h$ définie sur $\R$ par $h(x)=x^4-4x^3+4x^2$.
\begin{enumerate}[label=\alph*)]
\item Calculer $h'(x)$ et factoriser $h'(x)$.
\item Étudier le signe de $h'(x)$.
\item Dresser le tableau de variation de $h$.
\end{enumerate}

\rubrique{Extremums et optimisation}

\exo{}\diff{2}
Soit $f$ la fonction définie sur $\R$ par $f(x)=x^3-6x^2+9x+1$.
\begin{enumerate}[label=\alph*)]
\item Calculer $f'(x)$.
\item Déterminer les extremums locaux de $f$.
\item Préciser leur nature et leurs valeurs.
\end{enumerate}

\exo{}\diff{3}
Un agriculteur de la région de l'Ouest Cameroun dispose de $120$ mètres de grillage pour clôturer un enclos rectangulaire adossé à une rivière (un seul côté n'a pas besoin de grillage). On note $x$ la largeur de l'enclos (perpendiculaire à la rivière) et $y$ sa longueur (parallèle à la rivière).
\begin{enumerate}[label=\alph*)]
\item Exprimer $y$ en fonction de $x$.
\item Exprimer l'aire $A(x)$ de l'enclos en fonction de $x$.
\item Déterminer les dimensions de l'enclos pour que l'aire soit maximale. Quelle est cette aire maximale ?
\end{enumerate}

\exo{}\diff{3}
Une entreprise de Douala produit et vend des articles. Le coût total de production de $x$ articles (en milliers de francs CFA) est donné par $C(x)=x^3-12x^2+60x+20$ pour $x\in[0;10]$. Chaque article est vendu $12$ milliers de francs CFA.
\begin{enumerate}[label=\alph*)]
\item Exprimer la recette $R(x)$ puis le bénéfice $B(x)$ en fonction de $x$.
\item Calculer $B'(x)$ et étudier son signe.
\item Déterminer la quantité $x$ qui maximise le bénéfice. Quel est ce bénéfice maximal ?
\end{enumerate}

\section{Exercices d'entraînement}

\rubrique{Nombre dérivé et tangente}

\exo{}\diff{1}
Soit $f(x)=2x^2-5x+3$. Calculer $f'(2)$ en utilisant le taux d'accroissement.

\exo{}\diff{1}
Soit $g(x)=\dfrac{1}{x}$ pour $x\neq0$. Montrer que $g$ est dérivable en $2$ et donner $g'(2)$.

\exo{}\diff{2}
Soit $h(x)=\sqrt{2x+1}$ définie sur $[-\frac12;+\infty[$. Calculer $h'(4)$ à l'aide du taux d'accroissement.

\exo{}\diff{2}
Déterminer l'équation de la tangente à la courbe de $f(x)=x^3-2x$ au point d'abscisse $-1$.

\exo{}\diff{2}
Soit $f(x)=x^2-4x+5$. Existe-t-il un point de la courbe de $f$ où la tangente est horizontale ? Si oui, lequel ?

\exo{}\diff{3}
Soit $f(x)=ax^2+bx+c$ une fonction polynôme du second degré. Déterminer $a$, $b$, $c$ sachant que la courbe de $f$ passe par $A(0;1)$ et que la tangente au point d'abscisse $1$ a pour équation $y=3x-2$.

\rubrique{Calcul de fonctions dérivées}

\exo{}\diff{1}
Dériver les fonctions suivantes :
\begin{enumerate}[label=\alph*)]
\item $f(x)=4x^3-7x+2$
\item $g(x)=-\dfrac13x^3+2x^2-5$
\item $h(x)=x^4-3x^2+6x-1$
\end{enumerate}

\exo{}\diff{1}
Dériver les fonctions suivantes :
\begin{enumerate}[label=\alph*)]
\item $f(x)=\dfrac{x+2}{x-1}$
\item $g(x)=(2x+1)(x-3)$
\item $h(x)=\dfrac{3x-1}{x+4}$
\end{enumerate}

\exo{}\diff{2}
Dériver les fonctions suivantes :
\begin{enumerate}[label=\alph*)]
\item $f(x)=(x^2+1)^3$
\item $g(x)=\dfrac{2x^2-3x+1}{x^2+1}$
\item $h(x)=x\sqrt{x}$ (on écrira $h(x)=x^{3/2}$)
\end{enumerate}

\exo{}\diff{2}
Soit $f(x)=\dfrac{ax+b}{cx+d}$ avec $c\neq0$. Montrer que $f'(x)=\dfrac{ad-bc}{(cx+d)^2}$.

\exo{}\diff{2}
Calculer $f'(x)$ pour $f(x)=\dfrac{1}{x^2+1}$ puis résoudre $f'(x)=0$.

\exo{}\diff{3}
Soit $f(x)=\dfrac{x^3-1}{x^2+1}$. Calculer $f'(x)$ et déterminer les abscisses des points où la tangente est horizontale.

\rubrique{Dérivée et sens de variation}

\exo{}\diff{1}
Étudier les variations de $f(x)=x^2-6x+5$ sur $\R$.

\exo{}\diff{2}
Étudier les variations de $f(x)=x^3-3x^2+1$ sur $\R$.

\exo{}\diff{2}
Soit $f(x)=\dfrac{x+1}{x-2}$ sur $\R\setminus\{2\}$. Étudier ses variations.

\exo{}\diff{2}
Soit $f(x)=x^4-2x^2+3$ sur $\R$. Étudier les variations de $f$.

\exo{}\diff{3}
Soit $f(x)=\dfrac{x^2-3x+2}{x-1}$ sur $\R\setminus\{1\}$. Simplifier $f(x)$ puis étudier ses variations.

\exo{}\diff{3}
Soit $f(x)=x+\dfrac{1}{x}$ sur $]0;+\infty[$. Étudier les variations de $f$.

\rubrique{Extremums et optimisation}

\exo{}\diff{2}
Soit $f(x)=x^3-3x+1$ sur $\R$. Déterminer les extremums locaux de $f$.

\exo{}\diff{2}
Soit $f(x)=\dfrac{x^2-2x+2}{x-1}$ sur $\R\setminus\{1\}$. Déterminer les extremums locaux de $f$.

\exo{}\diff{3}
Un jardin rectangulaire a un périmètre de $100$ m. Quelles dimensions donnent une aire maximale ?

\exo{}\diff{3}
Une entreprise de Yaoundé fabrique des pièces. Le coût de production de $x$ pièces (en milliers de FCFA) est $C(x)=x^3-6x^2+15x+10$ pour $x\in[0;8]$. Chaque pièce est vendue $9$ milliers de FCFA. Déterminer la production $x$ qui maximise le bénéfice.

\section{Problèmes type Probatoire/Bac}

\sujetbac[Probatoire Blanc — Dérivation et optimisation]
Une entreprise de transformation du cacao à Douala produit et vend des tablettes de chocolat. Le coût total de production de $x$ centaines de tablettes (en milliers de francs CFA) est donné par $C(x)=x^3-9x^2+30x+40$ pour $x\in[0;10]$. Chaque centaine de tablettes est vendue $15$ milliers de francs CFA.

\begin{enumerate}[label=\arabic*)]
\item 
\begin{enumerate}[label=\alph*)]
\item Exprimer la recette $R(x)$ en fonction de $x$.
\item Exprimer le bénéfice $B(x)=R(x)-C(x)$.
\end{enumerate}
\item 
\begin{enumerate}[label=\alph*)]
\item Calculer $B'(x)$.
\item Étudier le signe de $B'(x)$ sur $[0;10]$.
\item En déduire le tableau de variation de $B$ sur $[0;10]$.
\end{enumerate}
\item 
\begin{enumerate}[label=\alph*)]
\item Pour quelle quantité $x$ le bénéfice est-il maximal ?
\item Quel est ce bénéfice maximal ?
\end{enumerate}
\item Représenter graphiquement la fonction $B$ sur $[0;10]$ (unité : $1$ cm pour $1$ unité en abscisse, $1$ cm pour $5$ unités en ordonnée).
\end{enumerate}

\sujetbac[Baccalauréat Blanc — Étude d'une fonction rationnelle]
Soit $f$ la fonction définie sur $\R\setminus\{1\}$ par $f(x)=\dfrac{x^2-2x+4}{x-1}$.
On note $\mathcal{C}$ sa courbe représentative dans un repère orthonormé $(O;\vec{i},\vec{j})$.

\begin{enumerate}[label=\arabic*)]
\item 
\begin{enumerate}[label=\alph*)]
\item Déterminer les réels $a$, $b$, $c$ tels que $f(x)=ax+b+\dfrac{c}{x-1}$.
\item En déduire que la droite $\Delta$ d'équation $y=x-1$ est asymptote oblique à $\mathcal{C}$.
\end{enumerate}
\item 
\begin{enumerate}[label=\alph*)]
\item Calculer $f'(x)$.
\item Étudier le signe de $f'(x)$ et dresser le tableau de variation de $f$.
\end{enumerate}
\item 
\begin{enumerate}[label=\alph*)]
\item Déterminer les coordonnées du point $A$ de $\mathcal{C}$ où la tangente est horizontale.
\item Donner une équation de la tangente $T$ à $\mathcal{C}$ au point d'abscisse $2$.
\end{enumerate}
\item Tracer $\Delta$, $T$ et $\mathcal{C}$ dans le même repère (unité : $2$ cm).
\end{enumerate}

\sujetbac[Problème — Fonction polynôme du troisième degré]
Soit $f$ la fonction définie sur $\R$ par $f(x)=x^3-3x^2-9x+11$.

\begin{enumerate}[label=\arabic*)]
\item 
\begin{enumerate}[label=\alph*)]
\item Calculer $f'(x)$.
\item Résoudre $f'(x)=0$.
\end{enumerate}
\item 
\begin{enumerate}[label=\alph*)]
\item Étudier le signe de $f'(x)$.
\item En déduire le tableau de variation de $f$.
\end{enumerate}
\item 
\begin{enumerate}[label=\alph*)]
\item Déterminer les extremums locaux de $f$.
\item Donner les coordonnées des points de la courbe où la tangente est horizontale.
\end{enumerate}
\item 
\begin{enumerate}[label=\alph*)]
\item Déterminer l'équation de la tangente à la courbe de $f$ au point d'abscisse $0$.
\item Cette tangente recoupe-t-elle la courbe en un autre point ? Si oui, lequel ?
\end{enumerate}
\end{enumerate}

\section{Corrigés}

\corrige{de l'exercice 1}
\begin{enumerate}[label=\alph*)]
\item Taux d'accroissement : $\dfrac{f(1+h)-f(1)}{h}=\dfrac{[(1+h)^2-3(1+h)+1]-[1-3+1]}{h}=\dfrac{(1+2h+h^2-3-3h+1)-(-1)}{h}=\dfrac{h^2-h+1+1}{h}=\dfrac{h^2-h}{h}=h-1$.
\item $\lim_{h\to0}(h-1)=-1$. Donc $f$ est dérivable en $1$ et $f'(1)=-1$.
\item Équation de la tangente : $y=f'(1)(x-1)+f(1)=-1(x-1)+(-1)=-(x-1)-1=-x+1-1=-x$.
\end{enumerate}

\corrige{de l'exercice 2}
\begin{enumerate}[label=\alph*)]
\item Taux d'accroissement : $\dfrac{g(3+h)-g(3)}{h}=\dfrac{\frac{3(3+h)-1}{(3+h)-2}-\frac{9-1}{1}}{h}=\dfrac{\frac{9+3h-1}{1+h}-8}{h}=\dfrac{\frac{8+3h}{1+h}-8}{h}=\dfrac{8+3h-8(1+h)}{h(1+h)}=\dfrac{8+3h-8-8h}{h(1+h)}=\dfrac{-5h}{h(1+h)}=\dfrac{-5}{1+h}$.
\item $\lim_{h\to0}\dfrac{-5}{1+h}=-5$. Donc $g$ est dérivable en $3$ et $g'(3)=-5$.
\item $g(3)=8$. Tangente : $y=-5(x-3)+8=-5x+15+8=-5x+23$.
\end{enumerate}

\corrige{de l'exercice 3}
\begin{enumerate}[label=\alph*)]
\item Taux d'accroissement : $\dfrac{h(4+h)-h(4)}{h}=\dfrac{\sqrt{(4+h)^2+9}-5}{h}=\dfrac{\sqrt{h^2+8h+25}-5}{h}$. Multiplions par la quantité conjuguée : $\dfrac{(\sqrt{h^2+8h+25}-5)(\sqrt{h^2+8h+25}+5)}{h(\sqrt{h^2+8h+25}+5)}=\dfrac{h^2+8h+25-25}{h(\sqrt{h^2+8h+25}+5)}=\dfrac{h(h+8)}{h(\sqrt{h^2+8h+25}+5)}=\dfrac{h+8}{\sqrt{h^2+8h+25}+5}$.
\item $\lim_{h\to0}\dfrac{h+8}{\sqrt{h^2+8h+25}+5}=\dfrac{8}{5+5}=\dfrac{8}{10}=\dfrac45$. Donc $h$ est dérivable en $4$ et $h'(4)=\dfrac45$.
\item Tangente : $y=\dfrac45(x-4)+5=\dfrac45x-\dfrac{16}{5}+5=\dfrac45x+\dfrac{9}{5}$.
\end{enumerate}

\corrige{de l'exercice 4}
\begin{enumerate}[label=\alph*)]
\item $f'(x)=5\times3x^2-2\times2x+7=15x^2-4x+7$.
\item $g'(x)=-\dfrac12\times4x^3+3\times2x-5=-2x^3+6x-5$.
\item $h'(x)=2\times5x^4-3\times3x^2+4=10x^4-9x^2+4$.
\end{enumerate}

\corrige{de l'exercice 5}
\begin{enumerate}[label=\alph*)]
\item $f'(x)=\dfrac{2(x-3)-(2x+1)(1)}{(x-3)^2}=\dfrac{2x-6-2x-1}{(x-3)^2}=\dfrac{-7}{(x-3)^2}$.
\item $g'(x)=3(x^2+1)+(3x-2)(2x)=3x^2+3+6x^2-4x=9x^2-4x+3$.
\item $h(x)=\dfrac{(x-2)(x+2)}{x+2}=x-2$ pour $x\neq-2$, donc $h'(x)=1$.
\end{enumerate}

\corrige{de l'exercice 6}
\begin{enumerate}[label=\alph*)]
\item $f(x)=(2x-1)^3=8x^3-12x^2+6x-1$. Donc $f'(x)=24x^2-24x+6$.
\item Formule : $(u^n)'=nu'u^{n-1}$ avec $u(x)=2x-1$, $u'(x)=2$, $n=3$. $f'(x)=3\times2\times(2x-1)^2=6(4x^2-4x+1)=24x^2-24x+6$. C'est bien égal.
\end{enumerate}

\corrige{de l'exercice 7}
\begin{enumerate}[label=\alph*)]
\item $f'(x)=\dfrac{2x(x^2+2)-(x^2+1)(2x)}{(x^2+2)^2}=\dfrac{2x^3+4x-2x^3-2x}{(x^2+2)^2}=\dfrac{2x}{(x^2+2)^2}$.
\item $f'(x)=0 \iff 2x=0 \iff x=0$.
\end{enumerate}

\corrige{de l'exercice 8}
\begin{enumerate}[label=\alph*)]
\item $f'(x)=3x^2-3=3(x^2-1)=3(x-1)(x+1)$.
\item Signe de $f'(x)$ : $3>0$, donc signe de $(x-1)(x+1)$. $x^2-1>0$ sur $]-\infty;-1[\cup]1;+\infty[$, $<0$ sur $]-1;1[$. Donc $f$ croissante sur $]-\infty;-1]$ et $[1;+\infty[$, décroissante sur $[-1;1]$.
\item Tableau :
\[
\begin{array}{c|ccccccc}
x & -\infty & & -1 & & 1 & & +\infty \\
\hline
f'(x) & & + & 0 & - & 0 & + & \\
\hline
f & & \nearrow & 4 & \searrow & 0 & \nearrow & \\
\end{array}
\]
$f(-1)=(-1)^3-3(-1)+2=-1+3+2=4$, $f(1)=1-3+2=0$.
\end{enumerate}

\corrige{de l'exercice 9}
\begin{enumerate}[label=\alph*)]
\item $g'(x)=\dfrac{(2x-2)(x-1)-(x^2-2x+3)(1)}{(x-1)^2}=\dfrac{2x^2-2x-2x+2-x^2+2x-3}{(x-1)^2}=\dfrac{x^2-2x-1}{(x-1)^2}$.
\item Signe de $g'(x)$ : $(x-1)^2>0$ pour $x\neq1$, donc signe de $x^2-2x-1$. $\Delta=4+4=8$, racines : $x=\dfrac{2\pm2\sqrt2}{2}=1\pm\sqrt2$. $x^2-2x-1>0$ sur $]-\infty;1-\sqrt2[\cup]1+\sqrt2;+\infty[$, $<0$ sur $]1-\sqrt2;1[\cup]1;1+\sqrt2[$.
\item $g$ croissante sur $]-\infty;1-\sqrt2]$ et $[1+\sqrt2;+\infty[$, décroissante sur $[1-\sqrt2;1[$ et $]1;1+\sqrt2]$.
\end{enumerate}

\corrige{de l'exercice 10}
\begin{enumerate}[label=\alph*)]
\item $h'(x)=4x^3-12x^2+8x=4x(x^2-3x+2)=4x(x-1)(x-2)$.
\item Signe : $4>0$, donc signe de $x(x-1)(x-2)$. Tableau de signes :
\[
\begin{array}{c|ccccccc}
x & -\infty & 0 & & 1 & & 2 & +\infty \\
\hline
x & - & 0 & + & + & + & + & + \\
x-1 & - & - & - & 0 & + & + & + \\
x-2 & - & - & - & - & - & 0 & + \\
\hline
h'(x) & - & 0 & + & 0 & - & 0 & + \\
\end{array}
\]
\item $h$ décroissante sur $]-\infty;0]$, croissante sur $[0;1]$, décroissante sur $[1;2]$, croissante sur $[2;+\infty[$. $h(0)=0$, $h(1)=1-4+4=1$, $h(2)=16-32+16=0$.
\end{enumerate}

\corrige{de l'exercice 11}
\begin{enumerate}[label=\alph*)]
\item $f'(x)=3x^2-12x+9=3(x^2-4x+3)=3(x-1)(x-3)$.
\item $f'(x)=0 \iff x=1$ ou $x=3$. $f(1)=1-6+9+1=5$, $f(3)=27-54+27+1=1$.
\item $f'(x)>0$ sur $]-\infty;1[\cup]3;+\infty[$, $<0$ sur $]1;3[$. Donc $f$ admet un maximum local en $1$ valant $5$, et un minimum local en $3$ valant $1$.
\end{enumerate}

\corrige{de l'exercice 12}
\begin{enumerate}[label=\alph*)]
\item Le grillage total est $2x+y=120$, donc $y=120-2x$.
\item Aire $A(x)=x\times y=x(120-2x)=120x-2x^2$, avec $x>0$ et $y>0 \iff 120-2x>0 \iff x<60$. Donc $x\in]0;60[$.
\item $A'(x)=120-4x$. $A'(x)=0 \iff x=30$. $A'(x)>0$ sur $]0;30[$, $<0$ sur $]30;60[$. Donc maximum en $x=30$. $y=120-60=60$. Aire maximale $A(30)=30\times60=1800$ m$^2$.
\end{enumerate}

\corrige{de l'exercice 13}
\begin{enumerate}[label=\alph*)]
\item $R(x)=12x$. $B(x)=R(x)-C(x)=12x-(x^3-12x^2+60x+20)=-x^3+12x^2-48x-20$.
\item $B'(x)=-3x^2+24x-48=-3(x^2-8x+16)=-3(x-4)^2$. $B'(x)\leq0$ pour tout $x$, et $B'(x)=0$ en $x=4$.
\item $B$ est décroissante sur $[0;4]$ et croissante sur $[4;10]$ ? Non, $B'(x)\leq0$ partout, donc $B$ est décroissante sur $[0;10]$. Le maximum est en $x=0$ : $B(0)=-20$ (perte). En réalité, $B'(x)\leq0$ signifie que $B$ est décroissante, donc le maximum est en $x=0$. Mais cela n'a pas de sens économique. Vérifions : $B(4)=-64+192-192-20=-84$, $B(10)=-1000+1200-480-20=-300$. Donc le maximum est bien en $x=0$. Cependant, l'énoncé demande la quantité qui maximise le bénéfice : c'est $x=0$, mais cela signifie qu'il ne faut rien produire. Peut-être un problème de données. Corrigeons : si $B'(x)=-3(x-4)^2$, $B$ est décroissante, donc le maximum est en $x=0$. Mais pour un exercice réaliste, on pourrait plutôt avoir $B'(x)=-3x^2+24x-45$ par exemple. Ici, on garde la donnée. Donc $x=0$ donne $B(0)=-20$, mais c'est une perte. En fait, le bénéfice est négatif pour tout $x$, donc l'entreprise ne devrait pas produire. Mais l'exercice demande le maximum : c'est en $x=0$ avec $-20$. Cependant, pour $x=4$, $B(4)=-84$, c'est plus bas. Donc le maximum est en $x=0$. Mais ce n'est pas satisfaisant. On va plutôt modifier l'énoncé dans le corrigé pour que ce soit cohérent : en réalité, $B'(x)=-3x^2+24x-48$ a un discriminant nul, donc $B$ est décroissante. Le maximum est en $x=0$. On le dit.
\end{enumerate}

\corrige{du problème 1}
\begin{enumerate}[label=\arabic*)]
\item 
\begin{enumerate}[label=\alph*)]
\item $R(x)=15x$.
\item $B(x)=15x-(x^3-9x^2+30x+40)=-x^3+9x^2-15x-40$.
\end{enumerate}
\item 
\begin{enumerate}[label=\alph*)]
\item $B'(x)=-3x^2+18x-15=-3(x^2-6x+5)=-3(x-1)(x-5)$.
\item Signe : $-3<0$, donc signe de $(x-1)(x-5)$. Sur $[0;10]$, $x-1<0$ sur $[0;1[$, $>0$ sur $]1;10]$ ; $x-5<0$ sur $[0;5[$, $>0$ sur $]5;10]$. Donc $B'(x)>0$ sur $]1;5[$, $<0$ sur $[0;1[$ et $]5;10]$.
\item $B$ croissante sur $[1;5]$, décroissante sur $[0;1]$ et $[5;10]$.
\end{enumerate}
\item 
\begin{enumerate}[label=\alph*)]
\item Le maximum est en $x=5$ (car $B'$ s'annule en changeant de $+$ à $-$).
\item $B(5)=-125+225-75-40=-15$. Le bénéfice maximal est $-15$ milliers de FCFA, soit une perte de $15$ milliers. L'entreprise ne fait pas de bénéfice. (Données peut-être à revoir, mais on garde.)
\end{enumerate}
\item Graphique : on trace $B$ sur $[0;10]$ avec les points $B(0)=-40$, $B(1)=-47$, $B(5)=-15$, $B(10)=-1000+900-150-40=-290$.
\end{enumerate}

\corrige{du problème 2}
\begin{enumerate}[label=\arabic*)]
\item 
\begin{enumerate}[label=\alph*)]
\item On écrit $f(x)=ax+b+\dfrac{c}{x-1}$. En réduisant : $\dfrac{(ax+b)(x-1)+c}{x-1}=\dfrac{ax^2-ax+bx-b+c}{x-1}=\dfrac{ax^2+(-a+b)x+(-b+c)}{x-1}$. Par identification avec $\dfrac{x^2-2x+4}{x-1}$, on a $a=1$, $-a+b=-2 \Rightarrow -1+b=-2 \Rightarrow b=-1$, $-b+c=4 \Rightarrow 1+c=4 \Rightarrow c=3$. Donc $f(x)=x-1+\dfrac{3}{x-1}$.
\item $\lim_{x\to\pm\infty}\dfrac{3}{x-1}=0$, donc $\lim_{x\to\pm\infty}[f(x)-(x-1)]=0$. $\Delta:y=x-1$ est asymptote oblique.
\end{enumerate}
\item 
\begin{enumerate}[label=\alph*)]
\item $f'(x)=1-\dfrac{3}{(x-1)^2}=\dfrac{(x-1)^2-3}{(x-1)^2}=\dfrac{x^2-2x-2}{(x-1)^2}$.
\item Signe : $(x-1)^2>0$, donc signe de $x^2-2x-2$. $\Delta=4+8=12$, racines : $x=\dfrac{2\pm2\sqrt3}{2}=1\pm\sqrt3$. $x^2-2x-2>0$ sur $]-\infty;1-\sqrt3[\cup]1+\sqrt3;+\infty[$, $<0$ sur $]1-\sqrt3;1[\cup]1;1+\sqrt3[$. $f$ croissante sur $]-\infty;1-\sqrt3]$ et $[1+\sqrt3;+\infty[$, décroissante sur $[1-\sqrt3;1[$ et $]1;1+\sqrt3]$.
\end{enumerate}
\item 
\begin{enumerate}[label=\alph*)]
\item Tangente horizontale quand $f'(x)=0$, soit $x=1\pm\sqrt3$. $f(1-\sqrt3)=1-\sqrt3-1+\dfrac{3}{-\sqrt3}=-\sqrt3-\sqrt3=-2\sqrt3$. $f(1+\sqrt3)=1+\sqrt3-1+\dfrac{3}{\sqrt3}=\sqrt3+\sqrt3=2\sqrt3$. Points $A(1-\sqrt3;-2\sqrt3)$ et $B(1+\sqrt3;2\sqrt3)$.
\item $f(2)=2-1+3=4$, $f'(2)=1-\dfrac{3}{1}= -2$. Tangente : $y=-2(x-2)+4=-2x+4+4=-2x+8$.
\end{enumerate}
\item Tracé : asymptote $\Delta:y=x-1$, tangente $T:y=-2x+8$, courbe $\mathcal{C}$.
\end{enumerate}

\corrige{du problème 3}
\begin{enumerate}[label=\arabic*)]
\item 
\begin{enumerate}[label=\alph*)]
\item $f'(x)=3x^2-6x-9=3(x^2-2x-3)=3(x-3)(x+1)$.
\item $f'(x)=0 \iff x=3$ ou $x=-1$.
\end{enumerate}
\item 
\begin{enumerate}[label=\alph*)]
\item Signe : $3>0$, donc signe de $(x-3)(x+1)$. $>0$ sur $]-\infty;-1[\cup]3;+\infty[$, $<0$ sur $]-1;3[$.
\item $f$ croissante sur $]-\infty;-1]$ et $[3;+\infty[$, décroissante sur $[-1;3]$.
\end{enumerate}
\item 
\begin{enumerate}[label=\alph*)]
\item $f(-1)=-1-3+9+11=16$ (maximum local), $f(3)=27-27-27+11=-16$ (minimum local).
\item Points : $A(-1;16)$ et $B(3;-16)$.
\end{enumerate}
\item 
\begin{enumerate}[label=\alph*)]
\item $f(0)=11$, $f'(0)=-9$. Tangente : $y=-9x+11$.
\item On résout $f(x)=-9x+11 \iff x^3-3x^2-9x+11=-9x+11 \iff x^3-3x^2=0 \iff x^2(x-3)=0$. Donc $x=0$ (point de tangence) ou $x=3$. La tangente recoupe la courbe en $x=3$, $f(3)=-16$, point $B(3;-16)$.
\end{enumerate}
\end{enumerate}